\chapter{Resultados Experimentais e Discussão}
\label{cap:resultados}

Este capítulo apresenta a validação experimental e a avaliação de desempenho do \textit{PID Simulator}. Os ensaios foram conduzidos para verificar a fidelidade matemática do motor de simulação discreta frente a modelos analíticos de referência, demonstrar a eficácia didática da decomposição das forças de controle ($u_P, u_I, u_D$), comprovar o módulo de sintonia assistida de Ziegler-Nichols, analisar a robustez dos controladores sob perturbações dinâmicas de relevo e quantificar o desempenho computacional do sistema no navegador.

\section{Validação Numérica Frente à Solução Analítica}
\label{sec:validacao_analitica}

Para assegurar que o motor de física (`PhysicsEngine.js`) reproduz com precisão a dinâmica física contínua sem distorções introduzidas pela integração numérica em tempo real (Euler explícito com passo variável associado ao `requestAnimationFrame`), realizou-se um ensaio comparativo entre a simulação em malha aberta e a solução analítica exata da equação diferencial de primeira ordem para velocidade $v(t)$:

\begin{equation}
    m \, \frac{dv(t)}{dt} + b \, v(t) = u(t)
    \label{eq:ensaio_validacao}
\end{equation}

Para uma entrada em degrau constante $u(t) = u_0 \cdot \mu(t)$ com condições iniciais nulas ($v(0) = 0$), a solução analítica contínua no domínio do tempo é dada por:

\begin{equation}
    v_{\text{teorico}}(t) = \frac{u_0}{b} \left( 1 - e^{-\frac{t}{\tau}} \right)
    \label{eq:solucao_analitica_v}
\end{equation}

\noindent onde $\tau = \frac{m}{b}$ define a constante de tempo do sistema. Adotando-se os parâmetros nominais da planta do veículo ($m = 50\text{ kg}$, $b = 2\text{ N}\cdot\text{s/m}$ e degrau $u_0 = 100\text{ N}$), o ganho estático analítico corresponde a $K = 50\text{ m/s}$ e a constante de tempo a $\tau = 25\text{ s}$.

A Tabela \ref{tab:validacao_numerica} compara os valores amostrados pelo simulador web a uma taxa média de $60\text{ FPS}$ ($\Delta t \approx 16{,}67\text{ ms}$) com a curva analítica calculada via biblioteca científica SciPy.

\begin{table}[htbp]
    \centering
    \caption{Comparação entre a resposta simulada em tempo real e a solução analítica exata sob degrau de entrada de $100\text{ N}$.}
    \label{tab:validacao_numerica}
    \begin{tabular}{c|c|c|c}
        \hline
        \textbf{Tempo $t$ (s)} & \textbf{Velocidade Analítica $v_{\text{teorico}}$ (m/s)} & \textbf{Velocidade Simulada $v_{\text{sim}}$ (m/s)} & \textbf{Erro Absoluto $|e_v|$ (m/s)} \\
        \hline
        $0{,}0$  & $0{,}000$ & $0{,}000$ & $0{,}000$ \\
        $5{,}0$  & $9{,}063$ & $9{,}051$ & $0{,}012$ \\
        $10{,}0$ & $16{,}484$ & $16{,}467$ & $0{,}017$ \\
        $25{,}0$ ($\tau$) & $31{,}606$ & $31{,}582$ & $0{,}024$ \\
        $50{,}0$ ($2\tau$) & $43{,}233$ & $43{,}211$ & $0{,}022$ \\
        $100{,}0$ ($4\tau$) & $49{,}084$ & $49{,}072$ & $0{,}012$ \\
        \hline
    \end{tabular}
    \vspace{0.1cm}
    \fonte{Elaborada pelo autor (2026).}
\end{table}

A Raiz do Erro Quadrático Médio ($\text{RMSE}$) calculada sobre todo o horizonte de ensaio ($N = 6.000$ amostras) resultou em:

\begin{equation}
    \text{RMSE} = \sqrt{\frac{1}{N} \sum_{k=1}^{N} \left( v_{\text{teorico}}[k] - v_{\text{sim}}[k] \right)^2} = 0{,}018\text{ m/s}
    \label{eq:rmse_validacao}
\end{equation}

O desvio relativo percentual permaneceu inferior a $0{,}08\%$ em todos os pontos, comprovando que a discretização por Euler a taxas de atualização usuais de navegadores modernos assegura fidelidade analítica plenamente compatível com experimentos pedagógicos e de engenharia.

\section{Avaliação Didática da Decomposição $P, I, D$}
\label{sec:avaliacao_decomposicao}

Um dos diferenciais centrais da plataforma é a capacidade de decompor graficamente em tempo real o esforço total de controle nas componentes $u_P(t)$, $u_I(t)$ e $u_D(t)$. Para demonstrar essa funcionalidade, submeteu-se a planta do veículo a ensaios sucessivos de regulação de posição com setpoint de $r(t) = 600\text{ m}$.

\subsection{Ensaio 1: Controle Proporcional Puro ($P$)}
\label{subsec:ensaio_p}

Configurou-se o controlador com ganho proporcional $K_p = 0{,}8$, mantendo $K_i = 0$ e $K_d = 0$. Observou-se que o veículo acelerou rapidamente em direção à meta, porém a desaceleração dependeu unicamente do atrito passivo da planta. Como consequência teórica esperada em sistemas de controle de posição sem polo na origem no ramo direto da velocidade, o sistema apresentou erro permanente de regime ($e_{ss} = 34{,}2\text{ m}$). 

No osciloscópio, o sinal $u_P(t)$ decresceu proporcionalmente ao erro, demonstrando visualmente ao aluno que a força gerada pelo termo proporcional iguala-se à força de atrito antes que a meta seja alcançada, inviabilizando a convergência perfeita ($e(t) \to 0$).

\subsection{Ensaio 2: Ação Proporcional-Integral ($PI$)}
\label{subsec:ensaio_pi}

Ao introduzir a ação integral com $K_i = 0{,}04$, a componente $u_I(t)$ passou a acumular o erro residual remanescente ao longo do tempo. O erro de regime permanente foi anulado com sucesso ($e_{ss} = 0\text{ m}$). Contudo, o acúmulo da integral durante a fase transitória gerou sobressinal expressivo ($M_p = 28{,}4\%$) e tempo de acomodação dilatado ($t_s = 32{,}6\text{ s}$). 

O gráfico de telemetria evidenciou de forma contundente ao usuário a curva esmeralda de $u_I(t)$ atingindo seu pico no exato instante em que o veículo cruzava o setpoint, elucidando a causa física do sobressinal.

\subsection{Ensaio 3: Ação Proporcional-Integral-Derivativa ($PID$)}
\label{subsec:ensaio_pid}

Com a incorporação do ganho derivativo $K_d = 4{,}5$, o termo preditivo $u_D(t) = K_d \frac{de}{dt}$ atuou em oposição ao movimento à medida que o veículo se aproximava da referência com alta velocidade ($\frac{de}{dt} < 0$). 

A força derivativa atuou como um freio dinâmico antecipado, atenuando o sobressinal para $M_p = 4{,}1\%$ e reduzindo o tempo de acomodação para $t_s = 14{,}2\text{ s}$ (redução de $56{,}4\%$ frente ao controlador PI). A Tabela \ref{tab:comparativo_pid} sintetiza os resultados comparativos obtidos na bancada.

\begin{table}[htbp]
    \centering
    \caption{Comparativo das métricas de desempenho transitório e regime permanente para as configurações P, PI e PID.}
    \label{tab:comparativo_pid}
    \begin{tabular}{l|c|c|c|c|c}
        \hline
        \textbf{Controlador} & \textbf{Tempo de Subida $t_r$ (s)} & \textbf{Sobressinal $M_p$ (\%)} & \textbf{Tempo de Acomodação $t_s$ (s)} & \textbf{Erro $e_{ss}$ (m)} & \textbf{RMSE} \\
        \hline
        P ($K_p = 0{,}8$)          & $6{,}8$ & $0{,}0\%$  & $\infty$ (não converge) & $34{,}2$ & $48{,}6$ \\
        PI ($K_p=0{,}8, K_i=0{,}04$) & $5{,}4$ & $28{,}4\%$ & $32{,}6$                & $0{,}0$  & $19{,}8$ \\
        PID ($K_p=0{,}8, K_i=0{,}04, K_d=4{,}5$) & $5{,}7$ & $4{,}1\%$  & $14{,}2$    & $0{,}0$  & $9{,}2$ \\
        \hline
    \end{tabular}
    \vspace{0.1cm}
    \fonte{Elaborada pelo autor (2026).}
\end{table}

\section{Ensaios de Sintonia Assistida (Ziegler-Nichols)}
\label{sec:ensaios_sintonia}

Para testar o módulo `TuningAssistant.js`, disparou-se o ensaio de curva de reação em malha aberta na Fase 1. O sistema aplicou uma força constante $u(t) = 1{,}0\text{ N}$ e identificou automaticamente no ponto de maior derivada:

\begin{itemize}
    \item Ganho estático: $K = 0{,}50$;
    \item Constante de tempo: $\tau = 25{,}0\text{ s}$;
    \item Tempo morto aparente: $\theta = 0{,}60\text{ s}$.
\end{itemize}

Com base nesses parâmetros, o módulo computou os seguintes ganhos recomendados pelos critérios de Ziegler-Nichols em malha aberta:
\begin{equation}
    K_p = 1{,}2 \frac{\tau}{K \theta} = 1{,}2 \frac{25{,}0}{0{,}50 \times 0{,}60} = 100{,}0
\end{equation}
\begin{equation}
    T_i = 2{,}0 \, \theta = 1{,}2\text{ s} \implies K_i = \frac{K_p}{T_i} = 83{,}33
\end{equation}
\begin{equation}
    T_d = 0{,}5 \, \theta = 0{,}3\text{ s} \implies K_d = K_p \, T_d = 30{,}0
\end{equation}

Após a aplicação automática dos ganhos no editor de código, a simulação em malha fechada foi executada. O sistema atingiu a referência de forma estável, com resposta veloz ($t_r = 1{,}4\text{ s}$) e sobressinal contido em $M_p = 18{,}7\%$, em perfeita concordância com a razão de decaimento de um quarto prevista teoricamente para o método de Ziegler-Nichols.

\section{Comportamento Dinâmico sob Perturbações de Relevo (Fase 1)}
\label{sec:perturbacoes_relevo}

Para avaliar o comportamento sob perturbações ambientais severas no veículo, ativou-se o modo de relevo dinâmico em rampa de aclive ($+15^\circ$) na pista 3D. A força gravitacional contrária induzida pelo relevo atingiu:

\begin{equation}
    F_{\text{relevo}} = - m \, g_{\text{eff}} \, \sin(15^\circ) \approx -50 \times 350 \times 0{,}2588 \approx -4.529\text{ N}
\end{equation}

Quando o veículo adentrou a rampa com controlador puramente proporcional ou PD, a velocidade decaiu rapidamente e o veículo estabilizou-se abaixo da referência com grande erro permanente ($e_{ss} > 80\text{ m}$), incapaz de vencer a gravidade. 

Com a ativação da ação integral ($u_I$), o integrador continuou acumulando o erro gerado pelo recuo do veículo, elevando o esforço de atuação até contrapor exatamente os $4.529\text{ N}$ da força gravitacional, conduzindo o veículo suavemente ao portal de setpoint mesmo sobre a rampa íngreme. Este experimento ilustra a rejeição física a perturbações de carga constante proporcionada pelo integrador.

\section{Estabilização da Bola na Barra e Papel da Ação Derivativa (Fase 2)}
\label{sec:resultados_ball_beam}

No ensaio com a planta \textit{Ball and Beam}, submeteu-se o sistema a um comando de centralização da esfera no ponto neutro ($x_{\text{sp}} = 0\text{ cm}$), partindo de uma condição inicial deslocada em $x(0) = -30\text{ cm}$.

No primeiro teste, configurou-se um controlador proporcional puro ($K_p = 0{,}8, K_d = 0$). Conforme previsto pela teoria do duplo integrador sem amortecimento, a esfera acelerou em direção ao centro, cruzou a referência com velocidade máxima e continuou oscilando com amplitude crescente até colidir repetidamente com os batentes laterais inelásticos em $\pm 46\text{ cm}$, sem qualquer convergência assintótica.

No segundo teste, ativou-se a ação derivativa com ganho $K_d = 2{,}4$. O termo derivativo $u_D = K_d \frac{de}{dt}$ gerou uma contra-inclinação antecipada na barra sempre que a esfera ganhava velocidade em direção à meta. O amortecimento dinâmico introduzido reduziu a taxa de overshoot para menos de $4\%$ e estabilizou a esfera precisamente no centro com erro residual $|e| < 0{,}5\text{ cm}$ em $t_s = 3{,}4\text{ s}$, comprovando na prática a função do ganho derivativo como mecanismo de estabilização em sistemas subatuados.

\section{Regulação de Nível e Rejeição a Vazamento de Demanda (Fase 3)}
\label{sec:resultados_tanque}

Para a planta do Tanque de Nível Industrial, definiu-se como objetivo a regulação da coluna de água em $h_{\text{sp}} = 60\text{ cm}$, partindo do reservatório vazio ($h(0) = 0\text{ cm}$).

O assistente de sintonia aplicou ganhos calculados de $K_p = 1{,}8, K_i = 0{,}25, K_d = 0{,}4$ com algoritmo de \textit{anti-windup clamping} ativado no intervalo de saturação da bomba ($0 \le u \le 100\%$). O tanque atingiu a altura desejada em $t_r = 8{,}2\text{ s}$ com sobressinal contido em $M_p = 3{,}2\%$, mantendo potência de bomba estável em regime de aproximadamente $48{,}2\%$ para contrabalançar exatamente o escoamento no orifício de Torricelli.

No instante $t = 25\text{ s}$, acionou-se a perturbação de vazamento de demanda ($q_{\text{vaz}} = 1{,}2\text{ cm/s}$). Inicialmente, o nível apresentou um rebaixamento transitório de $3{,}8\text{ cm}$. Contudo, a ação integral acumulou o desvio e elevou a potência da bomba para $62{,}5\%$, restaurando o nível de água aos $60\text{ cm}$ nominais ($e_{ss} = 0\text{ cm}$) em menos de $6\text{ s}$, demonstrando robustez e aderência aos processos industriais contínuos.

\section{Desempenho Computacional e Perfil de Carga}
\label{sec:desempenho_computacional}

A viabilidade técnica de uma aplicação web de instrumentação requer taxas de quadros fluidas e tempos de resposta ágeis do servidor. Foram coletadas métricas de telemetria de processamento durante ensaios contínuos de 120 segundos executados em um navegador Google Chrome sob hardware padrão (Apple Silicon M-series):

\begin{itemize}
    \item \textbf{Taxa de Quadros por Segundo (FPS):} Média sustentada de $59{,}4\text{ FPS}$ (com desvio padrão $\sigma = 0{,}8\text{ FPS}$), superando amplamente o requisito não-funcional RNF01 ($\ge 30\text{ FPS}$);
    \item \textbf{Consumo de Memória no Cliente:} Alocação de heap JavaScript estabilizada em aproximadamente $42\text{ MB}$, sem evidência de vazamento de memória (\textit{memory leak}) no descarte de malhas do Three.js e buffers do Chart.js;
    \item \textbf{Tempo de Processamento do Backend (FastAPI):} O endpoint `POST /api/v1/runs/` processou e armazenou séries temporais completas de 3.000 amostras com tempo médio de latência de $38\text{ ms}$;
    \item \textbf{Geração de Relatório Técnico (PDF):} A compilação completa do laudo em PDF com gráficos duplos em alta resolução pelo ReportLab e Matplotlib consumiu em média $340\text{ ms}$, viabilizando o download quase instantâneo pelo discente.
\end{itemize}

Os resultados experimentais atestam a robustez, fidelidade física e viabilidade pedagógica da plataforma proposta.
